\chapter{Circumcision}

\section*{Overview}
Circumcision is surgical removal of the foreskin covering the glans of the penis.

\section*{Why It Is Done}
It may be performed for cultural or personal reasons or to treat selected medical problems such as recurrent balanitis, scarring, or pathologic phimosis. It is not required for every person with a foreskin.

\section*{How It Is Performed}
After appropriate local or general anesthesia, the foreskin is separated, removed, and the edges are closed with absorbable sutures or allowed to heal according to the technique used.

\section*{Risks and Recovery}
Pain, bleeding, infection, swelling, excessive or inadequate skin removal, and altered sensation are possible. Keep the wound clean, follow dressing instructions, and avoid sexual activity until healing is complete.

\section*{Outcome and Follow-up}
The expected outcome depends on the reason for the procedure, the person's health, and whether additional treatment is needed. The clinician reviews the result with the patient, explains any next steps, and sets follow-up according to the findings and recovery.

\section*{When to Seek Medical Care}
Follow the procedure team's specific instructions. Seek urgent medical assessment for severe or worsening pain, uncontrolled bleeding, fever or signs of infection, trouble breathing, chest pain, fainting, new weakness or confusion, inability to urinate, or any rapidly worsening symptom. The appropriate warning signs depend on the procedure and the patient's medical condition.

\section*{References}
\begin{enumerate}
  \item MedlinePlus. Circumcision. U.S. National Library of Medicine; 2025.
  \item Current Medical Diagnosis and Treatment. McGraw-Hill; 2025.
\end{enumerate}
\clearpage
